% Generated draft: requires mathematical and source review.
\documentclass[11pt]{article}
\usepackage[margin=25mm]{geometry}
\usepackage{fontspec}
\setmainfont{Latin Modern Roman}
\usepackage{amsmath,amssymb,mathtools,unicode-math}
\setmathfont{Latin Modern Math}
\usepackage{enumitem,needspace,xurl,hyperref}
\setlength{\parindent}{0pt}
\setlength{\parskip}{5pt}
\setlength{\emergencystretch}{3em}
\newcommand{\N}{\mathbb{N}}
\newcommand{\Z}{\mathbb{Z}}
\newcommand{\Q}{\mathbb{Q}}
\newcommand{\R}{\mathbb{R}}
\newcommand{\C}{\mathbb{C}}
\newcommand{\F}{\mathbb{F}}
\newcommand{\A}{\mathbb{A}}
\newcommand{\PP}{\mathbb{P}}
\newcommand{\E}{\mathbb{E}}
\newcommand{\eps}{\varepsilon}
\newcommand{\dd}{\,\mathrm{d}}
\begin{document}
\section*{UnsolvedMath ID: 10400065}
\textbf{Source problem code:} catalogue entry
\subsection{Definitions and mathematical statement}
Let $(C_{\mathrm{graph}}^*,d)$ be the graph cochain complex intended by the source, and write $H_{\mathrm{graph}}^*=H^*(C_{\mathrm{graph}}^*,d)$. Graph-complex definitions depend on the admitted graphs, vertex valences, orientations/signs, relations, grading, differential, and completion; none is reproduced in the short item.

Find a natural geometric or conceptual object whose cohomology is this graph cohomology. A precise realization should identify the object and a natural comparison respecting the intended grading and other structure, rather than merely giving an abstract graded-vector-space coincidence. The source does not say whether the desired object is a topological space, a moduli object, a Lie algebra, or another cohomological construction.

The supplied summary attributes to Kontsevich an identification with cohomology of formal Hamiltonian vector fields. In standard versions of such correspondences, choices of graph complex and stable/relative/continuous Lie-algebra cohomology matter; the summary gives no such hypotheses, grading correspondence, or primary reference. Thus this is retained as a reported interpretation of a suitable version, not as an assertion that every graph cohomology is one unrestricted Lie-algebra cohomology or that the conceptual source question is completely answered.
\subsection{Short English statement}
Identify a natural object realizing the source's graph cohomology. Specify the graph complex and the comparison, and check the exact scope of the supplied Kontsevich/formal-Hamiltonian-vector-field interpretation.
\subsection{Sources}
\begin{itemize}
\item Ohtsuki, \emph{Problems on invariants of knots and 3-manifolds} (2002), Problem 3.15, PDF page 63, attributed to J. Roberts, \url{https://msp.org/gtm/2002/04/gtm-2002-04-024s.pdf}; accessed 2026-07-29, PDF-font-structure extraction.
\item Kontsevich, graph/Lie-algebra-cohomology correspondence attribution in the supplied summary, without a publication citation or exact stable/relative formulation.
\item Local registry \texttt{data/unsolvedmath/problems.json}, permanent UnsolvedMath ID $10400065$, source code catalogue entry, \emph{Problem 3.15}. Dataset attribution: Ulam AI and the UnsolvedMath Contributors.
\end{itemize}
\end{document}
